\section*{अध्याय \ref{ch:g10:habitable-planet} --- एक वासयोग्य ग्रह}

\begin{solution}{exo:g10:habitable-planet:1}
वह तापमान जिस पर आदर्श उत्सर्जक अंतरिक्ष में ठीक उतनी शक्ति विकिरित करता जितनी ग्रह
अपने तारे से अवशोषित करता है। पृथ्वी की सतह अधिक गरम है
(लगभग \qty{255}{K} के मुक़ाबले \qty{288}{K}) क्योंकि ग्रीनहाउस गैसें उसके द्वारा
उत्सर्जित अवरक्त का एक भाग अवशोषित करती हैं और उसका कुछ भाग वापस नीचे भेजती हैं।
\end{solution}

\begin{solution}{exo:g10:habitable-planet:2}
शुक्र: लगभग \qty{227}{K} के \omterm{def:g10:habitable-planet:equilibrium-temperature}{साम्य तापमान} के मुक़ाबले मापा गया \qty{737}{K},
यानी लगभग \qty{510}{K} की तापवृद्धि।
\end{solution}

\begin{solution}{exo:g10:habitable-planet:3}
$S = 1361/1.52^2 \approx \qty{589}{W/m^2}$ (तथ्य-तालिका सटीक माध्य दूरी के साथ
\qty{586}{W/m^2} देती है)।
\end{solution}

\begin{solution}{exo:g10:habitable-planet:4}
\qty{1}{bar}, \qty{330}{K}: द्रव (क्वथनांक, लगभग \qty{373}{K}, से नीचे)।
\qty{0.001}{bar}, \qty{250}{K}: बर्फ़ (दाब बर्फ़--वाष्प वक्र से ऊपर है, जो
\qty{250}{K} पर लगभग \qty{0.0008}{bar} पर है)।
\qty{100}{bar}, \qty{700}{K}: वाष्प --- तापमान क्रांतिक ताप से ऊपर है, इसलिए कोई
द्रव नहीं हो सकता।
\end{solution}

\begin{solution}{exo:g10:habitable-planet:5}
सूर्य के अतिरिक्त किसी अन्य तारे की परिक्रमा करने वाला ग्रह। पारगमन: ग्रह हर कक्षा में अपने
तारे के सामने से गुज़रता है और उसे थोड़ा धुँधला कर देता है। डगमगाहट: ग्रह का खिंचाव तारे
को हमारी दृष्टि रेखा के अनुदिश आगे-पीछे करता है, जिससे उसके प्रकाश का रंग आवर्ती रूप से
खिसकता है।
\end{solution}

\begin{solution}{exo:g10:habitable-planet:6}
$T = (1361 \times 0.706/(4 \times 5.67\times10^{-8}))^{1/4} \approx
\qty{255}{K} \approx \qty{-18}{\degreeCelsius}$। $A = 0$ के साथ:
$\approx \qty{278}{K} \approx \qty{5}{\degreeCelsius}$।
\end{solution}

\begin{solution}{exo:g10:habitable-planet:7}
$\sqrt 4 = 2$: $2 \times 0.99 = \qty{1.98}{au}$ से $2 \times 1.70 =
\qty{3.40}{au}$ तक। \qty{1}{au} पर ग्रह को पृथ्वी के किरणन का चार गुना मिलता:
भीतरी किनारे से बहुत अंदर, बहुत गरम।
\end{solution}

\begin{solution}{exo:g10:habitable-planet:8}
$S = 1361/0.723^2 \approx \qty{2604}{W/m^2}$;
$T_{\mathrm{eq}} = (2604 \times 0.23/(4\sigma))^{1/4} \approx \qty{227}{K}$।
सतह लगभग \qty{510}{K} अधिक गरम है: उसका घना कार्बन डाइऑक्साइड वायुमंडल
(\qty{92}{bar}) एक विशाल ग्रीनहाउस प्रभाव उत्पन्न करता है।
\end{solution}

\begin{solution}{exo:g10:habitable-planet:9}
\qty{6.36}{mbar} पर जल केवल लगभग \qty{0}{\degreeCelsius}
और \qty{0.5}{\degreeCelsius} के बीच द्रव रहता है। \qty{10}{\degreeCelsius} पर दाब
क्वथन वक्र से नीचे है: जल उबलकर वाष्प बन जाता है; रात में वह जम जाता है।
\end{solution}

\begin{solution}{exo:g10:habitable-planet:10}
$11.2/2.4 \approx 4.7$। पृथ्वी से पलायन के लिए गैस के अणु को चंद्रमा की तुलना में लगभग
पाँच गुना तेज़ चलना होता है; अरबों वर्षों में चंद्रमा के दुर्बल गुरुत्व ने उसकी गैसों को निकल
जाने दिया, जबकि पृथ्वी के गुरुत्व ने उनमें से अधिकांश को रोके रखा।
\end{solution}

\begin{solution}{exo:g10:habitable-planet:11}
$T = (1361 \times 0.50/(4\sigma))^{1/4} \approx \qty{234}{K}$, लगभग
\qty{-39}{\degreeCelsius}। बर्फ़ इतना अधिक सूर्य-प्रकाश परावर्तित करती कि ग्रह उसे बनाए
रखने लायक़ ठंडा बना रहता: एक स्व-प्रबलक चक्र।
\end{solution}

\begin{solution}{exo:g10:habitable-planet:12}
$T_{\mathrm{eq}} = (S_0(1-A)/(4\sigma d^2))^{1/4} = (S_0/(4\sigma))^{1/4}
(1-A)^{1/4} (d^2)^{-1/4}$; $(1361/(4\sigma))^{1/4} \approx \qty{278.3}{K}$ और
$(d^2)^{1/4} = \sqrt d$। चार गुना दूर: $T$ को $\sqrt 4 = 2$ से भाग।
\end{solution}

\begin{solution}{exo:g10:habitable-planet:13}
$\sqrt d = 278.3 \times 0.706^{1/4}/273 \approx 0.934$, इसलिए $d \approx
\qty{0.873}{au}$, भीतरी किनारे (\qty{0.99}{au}) के अंदर। \omterm{def:g10:habitable-planet:equilibrium-temperature}{साम्य
तापमान} ग्रीनहाउस प्रभाव की उपेक्षा करता है: वायुमंडल वाला वास्तविक ग्रह
$T_{\mathrm{eq}}$ से अधिक गरम होता है, इसलिए वह और बाहर ही, लगभग \qty{0.99}{au} पर,
अत्यधिक गरम होकर अपना जल खो देता है। किनारे पूरे वायुमंडल के मॉडलों से आते हैं,
केवल $T_{\mathrm{eq}}$ से नहीं।
\end{solution}

\begin{solution}{exo:g10:habitable-planet:14}
गलनांक, लगभग \qty{0}{\degreeCelsius}, और क्वथनांक,
लगभग \qty{7}{\degreeCelsius}, के बीच। \qty{15}{\degreeCelsius} पर झील
उबल जाती।
\end{solution}

\begin{solution}{exo:g10:habitable-planet:15}
$\sqrt{0.5} \approx 0.707$: $0.707 \times 0.99 \approx \qty{0.70}{au}$
से $0.707 \times 1.70 \approx \qty{1.20}{au}$ तक। ग्रह ठीक बाहरी किनारे पर है।
किनारे मॉडलों के परिणाम हैं और मापन सन्निकट है, इसलिए
यह कहना असंभव है कि वह भीतर है या बाहर।
\end{solution}

\begin{solution}{pb:g10:habitable-planet:1}
\textbf{1.} $0.00151/0.0485^2 \approx 0.642$।

\textbf{2.} $0.642 \times 1361 \approx \qty{874}{W/m^2}$।

\textbf{3.} दोनों के बीच: शुक्र से लगभग तीन गुना कम, मंगल से डेढ़ गुना
अधिक।

\textbf{4.} $365.25/11.18 \approx 32.7$।

\textbf{5.} $(874/(4 \times 5.67\times10^{-8}))^{1/4} \approx \qty{249}{K}$।

\textbf{6.} $(874 \times 0.23/(4\sigma))^{1/4} \approx \qty{173}{K}$।

\textbf{7.} कोई नहीं: \qty{249}{K} भी \qty{273}{K} से नीचे है। बादलों या बर्फ़ का चमकीला
आवरण \omterm{def:g10:habitable-planet:equilibrium-temperature}{साम्य तापमान} को
\qty{75}{K} से अधिक घटा देता है: ग्रह को क्या ढके हुए है, यह उतना ही मायने रखता है जितना उसकी दूरी।

\textbf{8.} $\sqrt{0.00151} \approx 0.0389$।

\textbf{9.} $0.99 \times 0.0389 \approx \qty{0.0385}{au}$।

\textbf{10.} $1.70 \times 0.0389 \approx \qty{0.066}{au}$।

\textbf{11.} हाँ: $0.0385 < 0.0485 < 0.066$।

\textbf{12.} $d_\odot = 0.0485/0.0389 \approx \qty{1.25}{au}$: सूर्य के \omterm{def:g10:habitable-planet:habitable-zone}{वासयोग्य क्षेत्र}
के भीतरी आधे भाग में, पृथ्वी और मंगल के बीच।

\textbf{13.} $1.055^{1/3} \approx 1.018$।

\textbf{14.} $11.2 \times \sqrt{1.055/1.018} \approx \qty{11.4}{km/s}$, लगभग
पृथ्वी जितना: उसका गुरुत्व वायुमंडल रोक सकता है (वायुमंडल है या नहीं, यह उसके इतिहास पर
निर्भर करता है)।

\textbf{15.} तारे की डगमगाहट उसकी गति का केवल वह भाग दिखाती है जो हमारी दृष्टि रेखा के
अनुदिश है; चूँकि कक्षा का झुकाव अज्ञात है, मिला द्रव्यमान वास्तविक द्रव्यमान तभी है जब कक्षा
किनारे से दिखे, अन्यथा वह न्यूनतम है।

\textbf{16.} त्रिक बिंदु का दाब, \qty{611.657}{Pa}, लगभग
\qty{0.006}{bar}।

\textbf{17.} $(874 \times 0.706/(4\sigma))^{1/4} \approx \qty{228}{K}$।

\textbf{18.} $288 - 228 = \qty{60}{K}$: प्रॉक्सिमा~b को \textbf{लगभग \qty{60}{K}} की
ग्रीनहाउस तापवृद्धि चाहिए, पृथ्वी के \qty{33}{K} की लगभग दुगुनी।
\end{solution}
